\documentclass[12pt,a4paper]{exam}
\usepackage[T1]{fontenc}
\usepackage[utf8]{inputenc}
\usepackage[provide=*,lithuanian]{babel}
\usepackage{lmodern}
\usepackage[left=1.5cm,right=1.5cm,top=2.4cm,bottom=1.5cm]{geometry}
\usepackage{amsmath,amssymb,tasks,array}
\pagestyle{headandfoot}
\runningheadrule
\firstpageheader{Mokytojo variantas}{\shortstack{Kontrolinis darbas:\\Daugianariai}}{Vilties variantas}
\runningheader{Daugianariai}{}{Vilties variantas}
\firstpagefooter{}{\thepage}{}
\runningfooter{}{\thepage}{}
\pointsdroppedatright
\bracketedpoints
\pointname{ tašk.}
\newcommand{\atsakymas}{\par\medskip\textbf{Atsakymas:}\ \framebox[9cm]{\rule{0pt}{0.7cm}}}
\begin{document}
\begin{center}\Large\textbf{Sprendimai ir vertinimo kriterijai}\end{center}
\section*{Taškų paskirstymas ir sunkumas}
Sunkumas vertinamas mokiniui, kuris jau mokėsi neigiamųjų laipsnių, daugianarių tapatybių ir pilnojo kvadrato išskyrimo.
\begin{center}\begin{tabular}{|c|c|p{10cm}|}\hline
Užduotis & Taškai & Sunkumas ir pagrindimas\\\hline
1 & 6 & Lengva--vidutinė: šeši atskiri klasifikavimo atvejai, svarbu atskirti kintamąjį nuo iracionaliosios konstantos.\\\hline
2 & 5 & Vidutinė--sunki: keli neigiamieji laipsniai, laipsnių dalyba ir vienanario sąvoka.\\\hline
3(a) & 7 & Sunki: ilgas suprastinimas, pilnojo kvadrato išskyrimas ir maksimumo pagrindimas.\\\hline
3(b) & 3 & Vidutinė: didelio rodiklio lyginumas ir parametrų panaudojimas.\\\hline
4 & 5 & Vidutinė: koeficientų palyginimas ir tapatybės patikra.\\\hline
5 & 4 & Vidutinė: trijų pilnųjų kvadratų išskyrimas ir minimumo pasiekiamumas.\\\hline
Iš viso & 30 & \\\hline
\end{tabular}\end{center}
\section*{1 užduotis (6 taškai)}
\begin{center}\begin{tabular}{|c|c|p{11cm}|}\hline
Dalis & Atsakymas & Paaiškinimas\\\hline
(a) & R, S & $25^{-7}$ ir $\sqrt6$ yra konstantos; kintamasis nėra vardiklyje ar po šaknimi.\\\hline
(b) & R, T & Vardiklyje yra kintamasis $x$; $x\ne0$.\\\hline
(c) & I & Kintamasis yra po kvadratine šaknimi; $x\leq5$.\\\hline
(d) & R, T & Vardiklyje yra $x-1$; $\sqrt\pi$ yra konstanta; $x\ne1$.\\\hline
(e) & R, T & $7x^{-2}-11=7/x^2-11$, todėl $x\ne0$.\\\hline
(f) & R, S & Vardiklis $\sqrt{11}$ yra nenulinė konstanta.\\\hline
\end{tabular}\end{center}
\textbf{Vertinimas:} po 1 tašką už kiekvieną tikslų raidžių rinkinį. Paaiškinimų ir apibrėžimo sričių mokinio sąlyga nereikalauja.
\section*{2 užduotis (5 taškai)}
Pradinio reiškinio apibrėžimo sąlygos: $x\ne0$, $y\ne0$.
\begin{align*}
(-3x^7y^{-5})^{12}&=3^{12}x^{84}y^{-60},\\
(27x^{-1}y^4)^{-20}&=3^{-60}x^{20}y^{-80},\\
(81x^4y^{-2})^{14}&=3^{56}x^{56}y^{-28}.
\end{align*}
Todėl
$$\frac{3^{12}x^{84}y^{-60}\cdot3^{-60}x^{20}y^{-80}}{3^{56}x^{56}y^{-28}}
=3^{-104}x^{48}y^{-112}=\boxed{\frac{x^{48}}{3^{104}y^{112}}}.$$
Pagal mokyklinę vienanario sampratą šis reiškinys \textbf{nėra vienanaris}, nes $y$ laipsnio rodiklis yra neigiamas. Vienanario kintamųjų laipsnių rodikliai turi būti neneigiami sveikieji skaičiai.
\par\textbf{Vertinimas:} po 1 tašką už kiekvieną iš trijų teisingai pakeltų laipsniu daugiklių (3 tšk.); 1 taškas už teisingą galutinį supaprastinimą; 1 taškas už išvadą ir jos pagrindimą. Apibrėžimo sritis atskirai nevertinama.
\section*{3 užduoties (a) dalis (7 taškai)}
\begin{align*}
(1-2x)^2(x+1)&=(1-4x+4x^2)(x+1)=4x^3-3x+1,\\
(x+1)^3&=x^3+3x^2+3x+1,\\
(x+2)(x-1)(x+3)&=(x^2+x-2)(x+3)=x^3+4x^2+x-6.
\end{align*}
Taigi
\begin{align*}
E(x)&=4x^3-3x+1-(x^3+3x^2+3x+1)\\
&\quad-3(x^3+4x^2+x-6)+14x^2-17\\
&=-x^2-9x+1\\
&=-\left(x^2+9x+\frac{81}{4}\right)+\frac{81}{4}+1\\
&=\boxed{-\left(x+\frac92\right)^2+\frac{85}{4}}.
\end{align*}
Vadinasi, $a=-1$, $b=9/2$, $c=85/4$.
Kadangi $(x+9/2)^2\geq0$, tai $E(x)\leq85/4$.
Lygybė pasiekiama, kai $x=-9/2$.
\par\textbf{Atsakymas:} didžiausia reikšmė yra $85/4=21{,}25$, ji gaunama, kai $x=-9/2$. Mažiausios reikšmės nėra, nes $E(x)\to-\infty$, kai $|x|\to\infty$.
\par\textbf{Vertinimas:} po 1 tašką už tris išskleidimus (3 tšk.); 1 taškas už $-x^2-9x+1$; 1 taškas už pilnojo kvadrato išskyrimą; 1 taškas už maksimumo reikšmę ir pagrindimą; 1 taškas už $x=-9/2$.
\section*{3 užduoties (b) dalis (3 taškai)}
Pažymėkime $N=1\overbrace{0\ldots0}^{666\text{ nulių}}2$. Šis skaičius yra lyginis, nes baigiasi skaitmeniu 2. Todėl
$$a^N=(-1)^N=1,\qquad 21{,}25:c=\frac{85}{4}:\frac{85}{4}=1.$$
Gauname
$$b(a^N+21{,}25:c)=\frac92(1+1)=\boxed9.$$
\textbf{Vertinimas:} 1 taškas už laipsnio reikšmę, remiantis rodiklio lyginumu; 1 taškas už dalmens reikšmę; 1 taškas už galutinį atsakymą.
\section*{4 užduotis (5 taškai)}
Išskleidžiame kairiąją pusę:
$$(ax-a)(2x^2-bx+c)=2ax^3-a(b+2)x^2+a(b+c)x-ac.$$
Palyginę atitinkamus koeficientus, gauname
$$\begin{cases}2a=-6,\\-a(b+2)=18,\\a(b+c)=-27,\\-ac=15.\end{cases}$$
Iš pirmos lygybės $a=-3$. Iš antros:
$$3(b+2)=18\quad\Longrightarrow\quad b=4.$$
Iš ketvirtos:
$$3c=15\quad\Longrightarrow\quad c=5.$$
Patikriname likusią lygybę: $a(b+c)=-3(4+5)=-27$. Visi koeficientai sutampa.
\par\textbf{Atsakymas:} $\boxed{a=-3,\ b=4,\ c=5}$.
\par\textbf{Vertinimas:} 1 taškas už bendrąjį išskleidimą; po 1 tašką už $a$, $b$, $c$ radimą (3 tšk.); 1 taškas už likusios koeficientų lygybės arba visos tapatybės patikrą.
\par\textit{Redakcinė pastaba:} pradinėje sąlygoje minėtas $d$ lygybėje nedalyvauja, todėl iš pataisytos sąlygos pašalintas. Pagal pradinį tekstą jis galėtų būti bet kuris realusis skaičius.
\section*{5 užduotis (4 taškai)}
\begin{align*}
9x^2-6x&=(3x-1)^2-1,\\
4y^2-12y&=(2y-3)^2-9,\\
z^2+8z&=(z+4)^2-16.
\end{align*}
Todėl
$$9x^2+4y^2+z^2-6x-12y+8z+27=(3x-1)^2+(2y-3)^2+(z+4)^2+1\geq1.$$
Visi trys kvadratai vienu metu lygūs nuliui, kai
$$x=\frac13,\qquad y=\frac32,\qquad z=-4.$$
\textbf{Atsakymas:} mažiausia reikšmė yra $\boxed1$.
\par\textbf{Vertinimas:} po 1 tašką už kiekvieną teisingai išskirtą pilnąjį kvadratą (3 tšk.); 1 taškas už minimumą ir jo pasiekiamumo pagrindimą.
\section*{Bendros vertinimo nuostatos}
Lygiaverčiai teisingi sprendimo būdai vertinami taip pat. Už tą pačią klaidą pakartotinai taškai neatimami: tolesni metodiškai teisingi žingsniai vertinami pagal mokinio gautą rezultatą, jei klaida iš esmės nesupaprastino užduoties. Tai taikoma ir 3(b) daliai, naudojant 3(a) dalyje gautus parametrus.
\par\medskip
\begin{center}
\textbf{Vertinimo lentelė (iš viso 30 taškų)}\par\medskip
\setlength{\tabcolsep}{5pt}
\renewcommand{\arraystretch}{1.3}
\begin{tabular}{|c|*{10}{c|}}
\hline
Taškai & 0--3 & 4--6 & 7--9 & 10--12 & 13--15 & 16--18 & 19--22 & 23--25 & 26--28 & 29--30\\\hline
Pažymys & 1 & 2 & 3 & 4 & 5 & 6 & 7 & 8 & 9 & 10\\\hline
\end{tabular}
\end{center}
Procentinis rezultatas $p=100t/30$, kur $t$ -- surinkti taškai.
Kai $p<65$, skaičius $p/10$ apvalinamas į didesnę pusę;
kai $p\geq65$, jis apvalinamas iki artimiausio sveikojo skaičiaus
(pusė apvalinama į didesnę pusę). Mažiausias pažymys -- 1.
Procentinis rezultatas prieš skaičiuojant pažymį neapvalinamas.

\end{document}